% Template LaTeX para TCC de graduação da UNIPAMPA — V2
% Copyright (c) 2026 Charles Quevedo Carpes.
% Autoria das contribuições próprias, organização e adaptações do template.
% This work may be distributed and/or modified under the conditions of the
% LaTeX Project Public License, version 1.3c or (at your option) any later version.
% License: https://www.latex-project.org/lppl/lppl-1-3c.txt (see LICENSE).
% This work has the LPPL maintenance status `author-maintained'.
% The Current Maintainer of this work is Charles Quevedo Carpes.
% Consulte AUTORIA.md para escopo, créditos e relação dos arquivos.

\chapter{INTRODUÇÃO}

A introdução apresenta o tema, a delimitação do estudo, o problema de pesquisa, a justificativa, os objetivos, uma síntese da metodologia e a organização do trabalho.

Este template adota o sistema autor-data. Uma citação indireta pode ser apresentada de modo narrativo, como em \citeonline{gil2022}, ou ao final da sentença \cite{gil2022}. A indicação de página em citações indiretas é opcional. Para maiores informações ou em caso de dúvidas deve ser  consultado o \textbf{Manual de Citações} disponível em: \url{https://sites.unipampa.edu.br/sisbi/files/2025/07/manual-de-normalizacao-de-citacao-sistema-de-chamada-e-notas-2025.pdf}


Para que as referências apareçam na seção correta, no final do texto, é necessário que elas tenham sido citadas em algum lugar do TCC. Vamos inserir uma frase aqui que cita todas as referências que temos no arquivo .bib que desse template. Isso fará o compilador inserir todas nas referências.

Estudos diversos abrangem áreas como literatura \cite{alves2000}, engenharia e tubulações \cite{bailona2006}, legislação federal \cite{brasil1988, brasil2024}, assessoria educacional \cite{cassol2003}, avaliação de revestimentos \cite{courtes2018}, auditoria bancária \cite{dantas2014}, jornalismo digital \cite{ferreira2003}, metodologia da pesquisa \cite{gil2022, ruiz2008}, teorias da personalidade \cite{hall2000}, redes de computadores \cite{jonack2006}, economia e globalização \cite{lacerda1998, mariano2005}, métodos educacionais \cite{capitulo2024}, tecnologias no ensino \cite{artigo2025}, pós-graduação \cite{dissertacao2023}, automação acadêmica \cite{evento2024}, normas técnicas \cite{abnt14724, unipampa2025manual}, biblioteconomia \cite{silva2012}, serviço social \cite{simoes2009}, softwares de apoio \cite{software2026}, diário oficial \cite{brasil1988}.

\section{Problema de pesquisa}

Apresente a pergunta que orienta a investigação de forma clara, delimitada e coerente com os objetivos.

\section{Objetivos}

\subsection{Objetivo geral}

Descreva o resultado principal que o trabalho pretende alcançar.

\subsection{Objetivos específicos}

Os objetivos específicos podem ser apresentados em alíneas:

\begin{alineas}
  \item identificar os elementos relevantes do problema;
  \item analisar os dados coletados;
  \item comparar os resultados obtidos; e
  \item apresentar as considerações finais.
\end{alineas}

\section{Organização do trabalho}

Explique brevemente o conteúdo das seções seguintes.
